% ======================================================================
% DODATAK B: OLAP ANALITIČKI UPITI
% ======================================================================
\chapter{OLAP analitički upiti}
\label{app:olap}

Ovaj prilog sadrži implementaciju tri reprezentativna OLAP analitička upita od ukupno pet korištenih u benchmarku. Odabrani su upiti koji ilustruju raspon SQL konstrukcija i time pokazuju zašto columnstore indeksi i in-memory pohrana imaju različit uticaj ovisno o tipu analitičke operacije. Za svaku od četiri konfiguracije postoje zasebne procedure -- \texttt{baseline.*}, \texttt{cs.*} i \texttt{im.*} su strukturno identične, ali upiti čitaju iz različitih storage slojeva (disk B-tree, disk columnstore, RAM). Iznimka je HTAP konfiguracija gdje su procedure registrirane u \texttt{htap.*} shemi ali čitaju iz \texttt{cs.*} sloja, što je ilustrirano u odjeljku~\ref{sec:b3}.

\section{Q1 -- Aktivnost postova po mjesecu (full-scan agregacija)}

Upit Q1 agregira sve postove po godini, mjesecu i tipu, računajući broj postova, prosječan score i ukupan broj pregleda. Skenira kompletnih $\sim$3,7 miliona redova tabele \texttt{Posts} bez filtera, što ga čini idealnim stresom za columnstore batch-mode izvršavanje. Na \texttt{baseline.*} konfiguraciji upit ide u row-mode i troši znatno više CPU ciklusa po redu; na \texttt{cs.*} i \texttt{htap.*} konfiguracijama columnstore indeks omogućava batch-mode obradu kolona, reducirajući broj instrukcija po redu za faktor 5--8.

\begin{lstlisting}[language=SQL, caption={baseline.usp\_OLAP\_Q1\_MonthlyActivity -- GROUP BY full-scan nad $\sim$3,7M redova}, label=lst:b1]
CREATE OR ALTER PROCEDURE baseline.usp_OLAP_Q1_MonthlyActivity AS
BEGIN
    SET NOCOUNT ON;
    SELECT
        YEAR(CreationDate)             AS Godina,
        MONTH(CreationDate)            AS Mjesec,
        PostTypeId,
        COUNT(*)                       AS BrojPostova,
        AVG(CAST(Score AS FLOAT))      AS ProsjekScore,
        SUM(CAST(ViewCount AS BIGINT)) AS UkupnoPregleda,
        SUM(AnswerCount)               AS UkupnoOdgovora
    FROM baseline.Posts
    GROUP BY YEAR(CreationDate), MONTH(CreationDate), PostTypeId
    ORDER BY Godina, Mjesec, PostTypeId;
END;
\end{lstlisting}

\section{Q3 -- Distribucija glasova po score bucketu (JOIN + CASE WHEN)}

Upit Q3 spaja tabele \texttt{Votes} ($\sim$10,1 miliona redova) i \texttt{Posts} ($\sim$3,7 miliona redova) putem \texttt{INNER JOIN}-a, grupira glasove prema score bucketu posta koristeći \texttt{CASE WHEN} logiku i broji glasove i jedinstvene postove po grupi. Ovakav JOIN s dvije velike tabele i naknadnom agregacijom predstavlja karakteristično BI opterećenje koje maksimalno koristi columnstore batch-mode koji može obraditi cijele stupce podataka u SIMD registrima, za razliku od row-mode koji obrađuje red po red.

\begin{lstlisting}[language=SQL, caption={baseline.usp\_OLAP\_Q3\_VoteDistribution -- JOIN Votes--Posts s CASE WHEN bucketingom}, label=lst:b2]
CREATE OR ALTER PROCEDURE baseline.usp_OLAP_Q3_VoteDistribution AS
BEGIN
    SET NOCOUNT ON;
    SELECT
        CASE
            WHEN p.Score < 0    THEN 'Negativan'
            WHEN p.Score = 0    THEN 'Nula'
            WHEN p.Score <= 5   THEN '1-5'
            WHEN p.Score <= 20  THEN '6-20'
            WHEN p.Score <= 100 THEN '21-100'
            ELSE                     '100+'
        END                      AS ScoreBucket,
        v.VoteTypeId,
        COUNT(*)                 AS BrojGlasova,
        COUNT(DISTINCT v.PostId) AS BrojPostova
    FROM baseline.Votes v
    INNER JOIN baseline.Posts p ON v.PostId = p.Id
    GROUP BY
        CASE WHEN p.Score < 0    THEN 'Negativan'
             WHEN p.Score = 0    THEN 'Nula'
             WHEN p.Score <= 5   THEN '1-5'
             WHEN p.Score <= 20  THEN '6-20'
             WHEN p.Score <= 100 THEN '21-100'
             ELSE '100+' END,
        v.VoteTypeId
    ORDER BY ScoreBucket, v.VoteTypeId;
END;
\end{lstlisting}

\section{Q5 -- Kohort aktivnost korisnika (CTE + dvostruka agregacija)}
\label{sec:b3}

Upit Q5 je najsloženiji u skupu: koristi CTE (\textit{Common Table Expression}) za izračun aktivnosti po korisniku i kvartalu, zatim agregira rezultate CTE-a u finalnu tabelu po kvartalima. Dvostruka agregacija (unutar CTE-a i nad CTE-om) tipičan je obrazac BI upita koji istovremeno testira sposobnost optimizatora da materijalizira privremene rezultate i izvrši naknadnu agregaciju. Upravo ovaj tip upita pokazuje najveći raspon performansi između konfiguracija: HTAP koji čita iz columnstore sloja (\texttt{cs.*}) izvršava batch-mode agregaciju u oba koraka, dok IM čita identičan upit iz RAM-a bez columnstore prednosti.

\begin{lstlisting}[language=SQL, caption={baseline.usp\_OLAP\_Q5\_UserCohortActivity -- CTE s dvostrukom agregacijom po kvartalu}, label=lst:b3]
CREATE OR ALTER PROCEDURE baseline.usp_OLAP_Q5_UserCohortActivity AS
BEGIN
    SET NOCOUNT ON;
    WITH UserQuarterActivity AS (
        SELECT
            OwnerUserId,
            YEAR(CreationDate)              AS Godina,
            DATEPART(QUARTER, CreationDate) AS Kvartal,
            COUNT(*)                        AS BrojPostova,
            SUM(Score)                      AS UkupniScore,
            SUM(CASE WHEN PostTypeId = 1 THEN 1 ELSE 0 END) AS BrojPitanja,
            SUM(CASE WHEN PostTypeId = 2 THEN 1 ELSE 0 END) AS BrojOdgovora
        FROM baseline.Posts
        WHERE OwnerUserId IS NOT NULL
        GROUP BY OwnerUserId, YEAR(CreationDate),
                 DATEPART(QUARTER, CreationDate)
    )
    SELECT
        Godina,
        Kvartal,
        COUNT(DISTINCT OwnerUserId)         AS AktivniKorisnici,
        SUM(BrojPostova)                    AS UkupnoPostova,
        AVG(CAST(BrojPostova AS FLOAT))     AS ProsjecnoPostovaPoKorisniku,
        SUM(UkupniScore)                    AS UkupniScore,
        SUM(BrojPitanja)                    AS UkupnoPitanja,
        SUM(BrojOdgovora)                   AS UkupnoOdgovora
    FROM UserQuarterActivity
    GROUP BY Godina, Kvartal
    ORDER BY Godina, Kvartal;
END;
\end{lstlisting}

\noindent\textit{Napomena o HTAP dual-store varijanti:} procedure Q1--Q5 za \texttt{htap.*} shemu su strukturno identične prikazanim, ali je svaka referenca na \texttt{baseline.*} tabelu zamijenjena odgovarajućom \texttt{cs.*} tablicom (\texttt{cs.Posts}, \texttt{cs.Users}, \texttt{cs.Votes}) -- čime OLAP workload čita iz disk-based columnstore sloja dok OLTP istovremeno piše u \texttt{htap.*} memorijsko-optimizirani sloj, ostvarujući fizičku izolaciju radnih tokova.
